% ch-ref-physics.tex -- reference: the Physics block and its nine sub-blocks
\chapter{Physics}
\label{ch:ref-physics}

\key{Physics} is required and must appear exactly once, but it carries no keys
of its own. It is a container for nine sub-blocks, one per physical effect,
each of which is optional and may appear at most once. Almost all of them are
switched with a single \key{Enable} key, and the term they control contributes
a coefficient to the flux in the height equation; when a term is disabled its
dimensionless group is set to zero and the term is not assembled.

This chapter lists the keys. What each term \emph{is} --- the flux it adds and
the assumptions behind it --- is in the theory part, which each section points
to.

\begin{note}[An absent \key{Physics} sub-block is a term switched off]
  A sub-block left out of the file reads as an empty one with its switch
  off: the term is not computed, and every other key takes the default listed
  in this chapter. For the four sub-blocks whose \key{Enable} defaults to
  \val{true} --- \key{LaplacePressure}, \key{SolutoCapillarity},
  \key{Evaporation} and \key{Source} --- leaving the block out is therefore
  not the same as writing it empty: an empty \code{[Physics.Evaporation]}
  switches evaporation on, a missing one leaves it off. A missing
  \key{ConcField} means \key{Model}~$=$~\val{"None"}. Every shipped case
  writes each sub-block it relies on.
\end{note}

\section{\key{Physics.ConcField}}
\label{sec:ref-physics:concfield}

This block decides how, or whether, the solute is transported, and its
\key{Model} key together with \key{Grid.Dimension} selects the whole solver
model. \val{"WM"} with \val{"xh"} gives the 1-D well-mixed model, \val{"WM"} with
\val{"xyh"} the 2-D one, \val{"CZ"} with \val{"xh"} the model that resolves
concentration through the thickness over a 1-D film, and \val{"CZ"} with
\val{"xyh"} the same over a 2-D film (the pseudo-3D case).
See \cref{ch:concentration}.

\begin{keytable}{Keys of \key{Physics.ConcField}\label{tab:physics-concfield-keys}}\key{Model} & string & \reqd & \val{"WM"} = well-mixed, depth-averaged concentration solved alongside the film; \val{"CZ"} = concentration resolved through the thickness, over a 1-D or a 2-D film (needs $\mathrm{Nz}\ge 2$); \val{"None"} = no concentration equation; \val{"CZ\_flow"} = accepted but not implemented. Required whenever the block is present; without the block the model is \val{"None"}. \\
\key{EnableConvection} & bool & \val{true} & Include the convective term in the concentration equation. Well-mixed models only. \\
\key{EnableDiffusion} & bool & \val{true} & Include the diffusive term in the concentration equation. Well-mixed models only. \\
\end{keytable}

\begin{warn}[\val{"CZ\_flow"} passes the range check and then does nothing]
  \val{"CZ\_flow"} is in the candidate list, so it is accepted without
  complaint, but no model implements it. The run falls back to the well-mixed
  model chosen by \key{Grid.Dimension} \emph{with the concentration equation
  switched off} --- the same thing \val{"None"} does, and the same thing an
  unknown string does if it somehow passes the check. The result is a
  plausible-looking run at constant concentration. Do not use it.
\end{warn}

When the concentration equation is not built --- \val{"None"}, \val{"CZ\_flow"}
or an unknown string --- \key{EnableConvection}, \key{EnableDiffusion},
\key{Fluid.Diffusivity} and the whole of \key{Numerics.ConcEquation} have no
effect.

\section{\key{Physics.LaplacePressure}}
\label{sec:ref-physics:laplace}

The surface-tension term: the tendency of the free surface to flatten itself.
It is what levels a film, what makes a drop a spherical cap, and what pulls
liquid into the dips of a rough substrate. It contains a fourth spatial
derivative, and it is the reason the height equation is fourth order, so
essentially every run keeps it on. The curvature is measured on the free
surface $h+s$, not on the film thickness, which is how the substrate shape
enters. See \cref{sec:terms:laplace}.

\begin{keytable}{Keys of \key{Physics.LaplacePressure}\label{tab:physics-laplace-keys}}\key{Enable} & bool & \val{true} & Switch on the capillary (Laplace) pressure term of the height equation. \\
\end{keytable}

\section{\key{Physics.SolutoCapillarity}}
\label{sec:ref-physics:soluto}

Solutal Marangoni flow: because surface tension depends on concentration, a
concentration gradient is a surface stress that drags liquid along the free
surface. It scales with $h^2$ rather than $h^3$, since it is a shear applied at
the surface rather than a body force, and it is a common origin of drying
streaks and coating mottle. See \cref{sec:terms:marangoni}.

\begin{keytable}{Keys of \key{Physics.SolutoCapillarity}\label{tab:physics-soluto-keys}}\key{Enable} & bool & \val{true} & Switch on solutal Marangoni flow. When false, $\Ma$ is forced to 0 and \key{Fluid.SurfaceTension.Coef} is silently ignored. \\
\end{keytable}

\begin{warn}[Disabling it also makes surface tension concentration-independent]
  $\Ma$ is not only the coefficient of the Marangoni stress; it is also what
  makes $\sigma$ depend on $c$ in the mobility of the capillary-pressure term.
  Setting \key{Enable = false} therefore removes the concentration dependence
  of surface tension \emph{entirely}, so that a drying film levels at the rate
  set by $\sigma_0$ alone. If you want the dependence without the Marangoni
  stress, this switch will not give it to you.
\end{warn}

\section{\key{Physics.DisjoiningPressure}}
\label{sec:ref-physics:disjoining}

The disjoining pressure
$\Pihat = (\Hstar/H)^n - (\Hstar/H)^m$ represents the interaction between
substrate and free surface once the film becomes very thin. It is what gives
the model partial wetting: a finite equilibrium contact angle, a moving contact
line, and a precursor film of thickness $\Hstar$ left outside the drop. Its
coefficient is fixed so that the angle you ask for is the one reproduced in the
lubrication limit, so changing the exponents or the precursor ratio does not
change the target angle --- it changes the thickness of the precursor and the
width of the transition region. Because $\Pihat \sim H^{-n}$, this is by a wide
margin the stiffest term in the model. See \cref{sec:terms:disjoining}.

\begin{keytable}{Keys of \key{Physics.DisjoiningPressure}\label{tab:physics-disjoining-keys}}\key{Enable} & bool & \val{false} & Switch on the disjoining pressure, i.e.\ partial wetting with a finite contact angle. Requires \key{Initial.SurfaceShape.PrecursorRatio} $>0$. \\
\key{ContactAngle} & double & \val{0} & Equilibrium contact angle $\thetae$ in \textbf{degrees}, range $[0:90]$, checked only when \key{Enable} is true. \val{0} gives zero disjoining pressure. \\
\key{ExponentN} & double & \val{3} & Exponent $n$ of the repulsive branch, non-dimensional. Must satisfy $n>m>1$, checked only when \key{Enable} is true. \\
\key{ExponentM} & double & \val{2} & Exponent $m$ of the attractive branch, non-dimensional, subject to the same constraint. \\
\end{keytable}

\begin{note}[The exponents default to $(n,m)=(3,2)$]
  Whether the block is absent or is written with \key{Enable = true} and no
  exponents, the pair is $(3,2)$.
\end{note}

\begin{warn}[A zero precursor ratio with disjoining pressure enabled aborts]
  $\Hstar$ is the thickness at which $\Pihat$ vanishes, so it cannot be zero.
  \key{Initial.SurfaceShape.PrecursorRatio = 0} together with
  \key{Enable = true} is fatal in \code{NonDimensionalNumbers::Build}, before
  any step is taken. The message in \file{FatalErrors.txt} reads
  \code{Physics.DisjoiningPressure requires
  Initial.SurfaceShape.PrecursorRatio > 0.}
  (\cref{sec:output:warnings}). Sizing the precursor is a modeling decision,
  not a formality: see \cref{sec:terms:disjoining}.
\end{warn}

\begin{note}[The \val{"CZ"} model ignores this block entirely]
  With \key{Physics.ConcField.Model = "CZ"} the height equation is built from
  capillary pressure, solutocapillarity, centrifugal force, vertical gravity
  and horizontal gravity only. \key{Physics.DisjoiningPressure} and the whole
  of \key{Physics.Source} have no effect, silently. Contact-line problems
  belong to the well-mixed models.
\end{note}

\section{\key{Physics.VerticalGravity}}
\label{sec:ref-physics:vgravity}

Gravity normal to the substrate: hydrostatic pressure drives liquid from high
parts of the free surface to low ones, so on a horizontal substrate it flattens
bumps and works alongside capillary pressure. Like capillary pressure it is
driven by $h+s$, so it too feels the substrate shape. If you are tilting the
substrate by an angle $\alpha$, the value to give here is $g\cos\alpha$. See
\cref{sec:terms:vgravity}.

\begin{keytable}{Keys of \key{Physics.VerticalGravity}\label{tab:physics-vgravity-keys}}\key{Enable} & bool & \val{false} & Switch on the hydrostatic, film-leveling gravity term. When false, $\Bo$ is forced to 0. \\
\key{Acceleration} & double & \val{9.80665} & Acceleration normal to the substrate in \si{\metre\per\second\squared}, non-negative; gives the Bond number $\Bo = \rho g L_0^2/\sigma_0$. Read even when \key{Enable} is false, and then discarded. \\
\end{keytable}

With the capillary length scale and \key{Acceleration = 9.80665} the Bond number is
exactly 1; \key{Acceleration} is the key to change when you want to move $\Bo$.

\section{\key{Physics.HorizontalGravity}}
\label{sec:ref-physics:hgravity}

An in-plane body force, that is, an inclined substrate. Unlike vertical
gravity the driving field is a constant vector independent of position, so
thicker parts of the film run faster and a front forms. For a substrate tilted
by $\alpha$, give the in-plane component $g\sin\alpha$ --- \num{4.903} at
$30^\circ$ --- and let its sign set the downhill direction. See
\cref{sec:terms:hgravity}.

\begin{keytable}{Keys of \key{Physics.HorizontalGravity}\label{tab:physics-hgravity-keys}}\key{Enable} & bool & \val{false} & Switch on the in-plane body force. When false, $\Pgx$ and $\Pgy$ are forced to 0. \\
\key{Acceleration} & array & \val{[0, 0]} & In-plane acceleration $(g_x,g_y)$ in \si{\metre\per\second\squared}. Length must be exactly~2, in 1-D runs as well. \\
\end{keytable}

\section{\key{Physics.Centrifugal}}
\label{sec:ref-physics:centrifugal}

The centrifugal term of spin coating: a rotating substrate throws the film
outwards, with a driving field proportional to the distance $r$ from the axis,
so the effect is strongest at the rim. In a 1-D (\val{"xh"}) run the liquid
thrown into the unsolved $y$ direction is accounted for as outflow, and
counted in \code{OutflowedVol} of \file{History.dat}. A rotating film is not
periodic, and a \val{"Periodic"} edge is refused with this term (except under
\key{Numerics.ManufacturedSolution}). See \cref{sec:terms:centrifugal}.

\begin{keytable}{Keys of \key{Physics.Centrifugal}\label{tab:physics-centrifugal-keys}}\key{Enable} & bool & \val{false} & Switch on the centrifugal (spin-coating) term. When false, $\Spn$ is forced to 0. \\
\key{RotationRate} & double & \val{0} & Angular velocity $\omega$ of the substrate in \si{\radian\per\second}, any value; $\Spn = \rho\omega^2 L_0^4/(\sigma_0 h_0)$. \\
\key{RotationAxis} & array & \val{[0, 0]} & In-plane position $(x,y)$ of the rotation axis, in meters. Length must be exactly~2. \\
\end{keytable}

\section{\key{Physics.Slip}}
\label{sec:ref-physics:slip}

A Navier slip length $\lambda$ at the substrate: the fluid slides over it with
the velocity $\lambda\,\partial u/\partial z$. Every term driven by a pressure
or a body force then has the mobility $h^{3}/3 + \lambda h^{2}$ in place of
$h^{3}/3$, and the Marangoni term $h^{2}/2 + \lambda h$ in place of $h^{2}/2$
(\cref{sec:terms:slip}). The dimensionless slip length is
$\Lambda = \lambda/h_0$, printed as \code{Lambda} under
\code{Nondimensional numbers} of the log when slip is on. Slip changes the
dynamics only: an equilibrium is the same with it. A missing block is no
slip; so is \key{Enable = true} with \key{Length = 0} (with a warning). The
\val{"CZ"} model refuses the block (\cref{ch:notverified}). With slip the
positivity of the film is guaranteed at one order less
(\cref{sec:disc:positivity-slip}).

\begin{keytable}{Keys of \key{Physics.Slip}\label{tab:physics-slip-keys}}\key{Enable} & bool & \val{false} & Switch on Navier slip at the substrate. When false, $\Lambda = 0$ and every term runs its no-slip expressions, bit for bit. \\
\key{Length} & double & \val{0.0} & Slip length $\lambda$ in meters, non-negative; $\Lambda = \lambda/h_0$. \\
\end{keytable}

\section{\key{Physics.Evaporation}}
\label{sec:ref-physics:evaporation}

Evaporation removes solvent, so the film thins and what is left grows more
concentrated. The rate is a product of a reference velocity, a function of
concentration, a function of position and a thin-film correction; the
concentration equation receives the same rate divided by the local thickness,
which is why thin regions concentrate fastest. See
\cref{sec:terms:evaporation}.

\begin{keytable}{Keys of \key{Physics.Evaporation}\label{tab:physics-evap-keys}}\key{Enable} & bool & \val{true} & Switch on solvent evaporation. When false the evaporation model is not built at all and the rate is 0 everywhere. \\
\key{Coef} & double & \val{1} & Reference evaporation velocity $E_0$ in \si{\metre\per\second}, non-negative; gives $\Ev = E_0 T_0/h_0$ and $\Pen = E_0 h_0/D_0$. \\
\key{SpatialFunc} & string & \val{"uniform"} & Spatial distribution of the flux: \val{"uniform"}, or \val{"negPow"}, the radial power $k(1-r/R')^{-n}$ of \cref{eq:evap-negpow}, which for $n>0$ diverges as a negative power toward the rim of the pool (capped at 10), blended over 5 cells into a quarter cosine that takes over the curve's value and slope, peaks at \key{NegPowRmax} and is constant beyond it. \\
\key{ConcFunc} & string & \val{"Raoult"} & Concentration dependence: \val{"Raoult"} for Raoult's law using the molar-mass ratio, or \val{"PowerN"} for $\max(1-c-\mathrm{ConcR},0)^{\mathrm{PowerN}}$. \\
\key{NegPowExponent} & double & \val{0.25} & Exponent $n$ of the \val{"negPow"} distribution $k(1-r/R')^{-n}$, non-dimensional, less than 1: positive gives a rate diverging at the rim, 0 a uniform one, negative one falling to zero there; the prefactor $k=\frac{1}{2}(1-n)(2-n)$ keeps the mean over the pool at 1. Used only when \key{SpatialFunc} is \val{"negPow"}. \\
\key{NegPowRmax} & double & \val{1} & Radius $R$ of the pool of the \val{"negPow"} distribution in meters, positive, measured from the origin of the grid coordinates. Used only with \val{"negPow"}. \\
\key{PowerN} & double & \val{0.5} & Exponent of the \val{"PowerN"} law, non-dimensional, any value; with 0 the rate does not fall with $c$, and the solvent limit stops it at $c = 1$ (\cref{sec:terms:evaporation}). Used only with \val{"PowerN"}. \\
\key{ConcR} & double & \val{0.0} & Residual concentration offset in the \val{"PowerN"} law, mass fraction in $[0:1)$: evaporation stops at $c = 1-\mathrm{ConcR}$. Used only with \val{"PowerN"}. \\
\key{EndConcType} & string & \val{"mean"} & Which concentration statistic is compared against \key{EndConcValue}: \val{"mean"}, \val{"max"} or \val{"min"}; any other string is refused. \\
\key{EndConcValue} & double & \val{0.95} & Dry-out threshold, mass fraction in $(0:1]$. The run stops with \code{Dry out} as soon as the chosen statistic exceeds it. \\
\end{keytable}

\begin{warn}[The dry-out test runs even with evaporation disabled]
  \key{EndConcValue} is applied whether or not \key{Enable} is true, and a
  file that leaves \key{Physics.Evaporation} out gets the default
  \val{0.95}. A film whose concentration reaches the threshold by any other
  route therefore stops with \code{Dry out} even with evaporation off.
\end{warn}

\key{PowerN}, \key{ConcR},
\key{EndConcType} and \key{EndConcValue} are not exercised by any case in
\file{cases/}.

\section{\key{Physics.Source}}
\label{sec:ref-physics:source}

Liquid supplied to the film through a nozzle, optionally carried on a rotating
arm --- the model of a dispensing head. Each cell receives supply in
proportion to the fraction of it that the nozzle footprint covers, and that
fraction is measured by subdividing the cell. See \cref{sec:terms:source}.

\begin{keytable}{Keys of \key{Physics.Source}\label{tab:physics-source-keys}}\key{Enable} & bool & \val{true} & Switch on the source term. Note the default is \textbf{true}: writing the block without \key{Enable} turns the nozzle on; leaving the block out leaves it off. \\
\key{Radius} & double & \val{1} & Radius of the nozzle footprint in meters, positive. \\
\key{InletRate} & double & \val{0} & Steady-state supply velocity in \si{\metre\per\second} (volume flux per unit area), non-negative; non-dimensionalized by $h_0/T_0$. \\
\key{ConcType} & string & \val{"local"} & The concentration the supplied liquid arrives at: \val{"local"}, that of the film where it lands (it thickens the film without changing $c$); \val{"supplied"}, \key{Conc}, mixed through the thickness of the column it lands on (\cref{sec:terms:source}). Any other string is refused. \\
\key{Conc} & double & \val{0} & Concentration of the supplied liquid, mass fraction in $[0:1]$; outside it is refused with \key{ConcType = "supplied"}. Used only with \val{"supplied"}. \\
\key{TimeConstant} & double & \val{0} & Exponential rise time $\tau$ of the flow rate in seconds, non-negative: $\mathrm{rate} = \mathrm{InletRate}\,(1-e^{-t/\tau})$. \val{0} means full rate from the first step. \\
\key{ArmLength} & double & \val{1} & Length of the rotating arm carrying the nozzle, in meters, non-negative. \val{0} pins the nozzle at \key{ArmCenter}. \\
\key{ArmCenter} & array & \val{[0, 0]} & Position $(x,y)$ of the arm pivot, in meters. Length must be exactly~2. \\
\key{RotationRate} & double & \val{0} & Angular velocity of the arm in \si{\radian\per\second}, any value. Distinct from \key{Physics.Centrifugal.RotationRate}. \\
\key{StartAngle} & double & \val{0} & Arm angle at $t=0$, in radians, any value. \\
\key{RampTime} & double & \val{0} & Ramp-up and ramp-down time of the arm's rotation rate in seconds, non-negative. \\
\key{PlateauTime} & double & \val{1e10} & How long the rotation rate stays at its plateau before ramping down, in seconds, non-negative. \\
\key{SubRegions} & int & \val{1} & Number of sub-divisions per cell used when measuring the covered fraction, at least 1. Larger values smooth the source as the nozzle moves. \\
\end{keytable}

\begin{note}[\key{ConcType} decides whether \key{Conc} is used]
  By default (\key{ConcType = "local"}) the supplied liquid takes the
  concentration already present where it lands, and \key{Conc} is not used:
  setting it alone does not dispense liquid of a different composition. With
  \key{ConcType = "supplied"} the liquid arrives at \key{Conc}, and the
  concentration under the nozzle follows from mixing it into the film
  (\cref{eq:source-mixing}). \code{SuppliedResin} in \file{History.dat}
  counts the solute either way brings. Without a concentration equation
  (\key{Physics.ConcField.Model = "None"}) the key has no effect.
\end{note}
